
\section{Background}

\subsection{Singular Value Decomposition}

For any matrix $W \in \mathbb{R}^{m \times n}$, the SVD is:

\[
W = U \Sigma V^\top
\tag{3.1}
\]

where $U \in \mathbb{R}^{m \times m}$ and $V \in \mathbb{R}^{n \times n}$ are 
orthogonal, and $\Sigma \in \mathbb{R}^{m \times n}$ is diagonal with 
non-negative entries $\sigma_1 \geq \sigma_2 \geq \dots \geq 0$.

\textbf{Truncated SVD.} Keeping the top-$r$ singular values gives the best rank-$r$ 
approximation in Frobenius norm:

\[
W_r = U_{:, :r} \Sigma_{:, :r} V_{:r, :}^\top
\tag{3.2}
\]

with squared error $\|W - W_r\|_F^2 = \sum_{i > r} \sigma_i^2$, or equivalently $\|W - W_r\|_F = \sqrt{\sum_{i > r} \sigma_i^2}$.

\textbf{Symmetric factorization.} For numerical stability, it is common to split 
$\Sigma^{1/2}$ between the two factors:

\[
U_r = U_{:, :r} \Sigma_{:, :r}^{1/2}, \qquad
V_r = \Sigma_{:r, :}^{1/2} V_{:r, :}^\top
\tag{3.3}
\]

so that $W_r = U_r V_r$.

\subsection{The Fisher Information Matrix}

For a probabilistic model $p_\theta(y | x)$, the Fisher information is:

\[
\mathcal{F}_{ij} = \mathbb{E}_{x, y}\!\left[\frac{\partial \log p_\theta(y | x)}{\partial \theta_i} \cdot \frac{\partial \log p_\theta(y | x)}{\partial \theta_j}\right]
\tag{3.4}
\]

For a linear layer $y = Wx + b$, the Fisher with respect to $x$ is:

\[
\mathcal{F}_x = \mathbb{E}_{x}\!\left[\frac{\partial \mathcal{L}}{\partial x} \frac{\partial \mathcal{L}}{\partial x}^\top\right]
\tag{3.5}
\]

The \textbf{diagonal} of this Fisher measures the importance of each input 
dimension. We use a cheap proxy: since $\partial \mathcal{L} / \partial x$ is 
proportional to $W^\top \cdot \partial \mathcal{L} / \partial y$, and since 
we assume the downstream loss has roughly uniform sensitivity to $y$, the 
diagonal Fisher is approximately:

\[
\mathcal{F}_{x, i} \approx \mathbb{E}_{x}\!\left[x_i^2\right] = F_i
\tag{3.6}
\]

This is a heuristic approximation under the assumptions that (i) the downstream loss has roughly uniform sensitivity across output dimensions, and (ii) cross-channel correlations are negligible. It is not an exact identity; we use it as a cheap proxy for the true Fisher diagonal. This \textbf{activation energy} is cheap to compute and correlates well with the true Fisher in practice.

\subsection{Streaming Covariance Accumulation}

Given a stream of activation vectors $x_1, x_2, \dots, x_N \in \mathbb{R}^d$, 
we want to compute:

\[
C = \frac{1}{N} \sum_{t=1}^N x_t x_t^\top, \qquad
F = \frac{1}{N} \sum_{t=1}^N x_t^2
\tag{3.7}
\]

Naively this requires storing all $x_t$, which costs $\mathcal{O}(Nd)$ memory. 
The \textbf{streaming} formulation accumulates only the running sums:

\[
C^{(k)} = C^{(k-1)} + \sum_{t \in \text{batch}_k} x_t x_t^\top, \qquad
N^{(k)} = N^{(k-1)} + |\text{batch}_k|
\tag{3.8}
\]

with $\mathcal{O}(d^2)$ memory. We use this in A3.

\textbf{Padding handling.} Padding tokens (all-zeros or mask-tagged) are filtered 
out before accumulation.

\subsection{LayerNorm}

LayerNorm (Ba et al., 2016) normalizes activations across the feature 
dimension:

\[
\text{LN}(x) = \gamma \odot \frac{x - \mu}{\sqrt{\mathrm{Var}(x) + \epsilon}} + \beta, \qquad
\text{where } \mathrm{Var}(x) = \mathbb{E}[(x - \mu)^2], \quad \sigma = \sqrt{\mathrm{Var}(x) + \epsilon}
\tag{3.9}
\]

where $\gamma, \beta \in \mathbb{R}^d$ are learnable parameters.

\textbf{Distributional shift.} When a preceding layer is replaced by a factorized 
approximation, its output distribution changes. The LayerNorm that follows 
sees a different input distribution and produces different normalized 
activations. Recalibrating $\gamma$ can partially absorb this shift.

\subsection{LoRA Recovery}

LoRA (Hu et al., 2021) parameterizes weight updates as:

\[
\Delta W = \frac{\alpha}{r} B A
\tag{3.10}
\]

with $A \in \mathbb{R}^{r \times k}$, $B \in \mathbb{R}^{d \times r}$, 
$B = 0$ at initialization. Only $A, B$ are trainable.

\textbf{Interaction with A3.} We add LoRA adapters to the factorized layer, 
yielding:

\[
y = (x V^\top) U^\top + (x A^\top) B^\top \cdot s + b
\tag{3.11}
\]

where $s = \alpha / r$ is the scaling factor. The frozen $U, V$ provide the 
bulk of the computation; LoRA provides the recovery signal.

\subsection{Notation Summary}

\begin{table}[H]
\centering
\small
\begin{tabular}{lll}
\toprule
\textbf{Symbol} & \textbf{Meaning} & \textbf{Dimension} \\
\midrule
$W$ & Weight matrix & $d_\text{out} \times d_\text{in}$ \\
$U$ & Left factor & $d_\text{out} \times r$ \\
$V$ & Right factor & $r \times d_\text{in}$ \\
$\Sigma$ & Singular value diagonal & $r \times r$ \\
$r$ & Compression rank & scalar \\
$\tau$ & Error threshold & scalar \\
$F$ & Empirical Fisher diagonal & $d_\text{in}$ \\
$C_X$ & Activation covariance & $d_\text{in} \times d_\text{in}$ \\
$x$ & Input activation & $d_\text{in}$ \\
$y$ & Output activation & $d_\text{out}$ \\
$\sigma_i$ & $i$-th singular value & scalar \\
$\epsilon$ & Numerical stability & scalar ($10^{-6}$) \\
$\gamma$ & LayerNorm scale & $d_\text{model}$ \\
$\beta$ & LayerNorm shift & $d_\text{model}$ \\
$\sigma_\text{target}$ & Pre-compression std & scalar \\
$\sigma_\text{current}$ & Post-compression std & scalar \\
$N$ & Calibration samples & scalar \\
$L$ & Sequence length & scalar \\
$B$ & Batch size & scalar \\
\bottomrule
\end{tabular}
\caption{Mathematical symbols, operations, and tensor dimensions used throughout this paper.}
\label{tab:notation_summary}
\end{table}
